% =============================================================================
% CAPÍTULO 4: RESULTADOS E DISCUSSÃO
% =============================================================================

\chapter{Resultados e Discussão}

\section{Resultados das Simulações Numéricas}

\subsection{Validação do Modelo Numérico}

O modelo numérico desenvolvido foi validado através da comparação com resultados experimentais disponíveis na literatura. A Figura \ref{fig:validacao_modelo} apresenta a comparação entre as frequências naturais calculadas e os valores experimentais reportados por \cite{exemplo_referencia}.

\begin{figure}[H]
    \centering
    \includegraphics[width=0.7\textwidth]{figures/validacao_modelo.png}
    \caption{Comparação entre frequências naturais calculadas e experimentais}
    \label{fig:validacao_modelo}
\end{figure}

Os resultados mostram uma excelente concordância entre os valores simulados e experimentais, com erro relativo médio inferior a 3\%, validando a fidelidade do modelo numérico desenvolvido.

\subsection{Análise de Resposta Dinâmica}

A análise da resposta dinâmica do sistema em condições normais revelou características espectrais típicas de sistemas rotativos, incluindo picos nas frequências naturais e harmônicos da velocidade de rotação. A Figura \ref{fig:espectro_normal} apresenta o espectro de potência para o sistema em condição normal.

\begin{figure}[H]
    \centering
    \includegraphics[width=0.8\textwidth]{figures/espectro_normal.png}
    \caption{Espectro de potência do sistema em condição normal}
    \label{fig:espectro_normal}
\end{figure}

\section{Análise de Espectros de Ordem Superior}

\subsection{Bi-espectros para Diferentes Condições}

A análise de bi-espectros revelou padrões característicos para cada tipo de falha estudado. A Figura \ref{fig:bispectro_comparacao} apresenta uma comparação dos bi-espectros para as diferentes condições analisadas.

\begin{figure}[H]
    \centering
    \includegraphics[width=0.9\textwidth]{figures/bispectro_comparacao.png}
    \caption{Comparação de bi-espectros para diferentes condições}
    \label{fig:bispectro_comparacao}
\end{figure}

\subsubsection{Sistema Normal}

O bi-espectro do sistema em condição normal apresenta baixa magnitude e distribuição relativamente uniforme, indicando ausência de não linearidades significativas.

\subsubsection{Sistema com Trinca}

O sistema com trinca apresenta picos característicos no bi-espectro, localizados em frequências específicas relacionadas à severidade da falha. Estes picos são consistentes com os padrões reportados na literatura para sistemas com trincas.

\subsubsection{Sistema com Desalinhamento}

O desalinhamento resulta em padrões distintos no bi-espectro, com picos em frequências relacionadas às forças desbalanceadas e acoplamentos não lineares entre os modos de vibração.

\subsection{Características Extraídas}

A Tabela \ref{tab:caracteristicas_hos} apresenta um resumo das características extraídas dos bi-espectros para as diferentes condições.

\begin{table}[H]
    \centering
    \caption{Características HOS extraídas para diferentes condições}
    \begin{tabular}{lccc}
        \toprule
        Característica & Normal & Trinca & Desalinhamento \\
        \midrule
        Magnitude máxima & 0.12 & 0.45 & 0.38 \\
        Energia total & 1.23 & 4.67 & 3.89 \\
        Entropia & 2.34 & 1.87 & 2.01 \\
        Número de picos & 2 & 8 & 6 \\
        \bottomrule
    \end{tabular}
    \label{tab:caracteristicas_hos}
\end{table}

\section{Validação Física}

\subsection{Comparação com Literatura}

A validação física dos resultados foi realizada através da comparação com estudos experimentais reportados na literatura. A Figura \ref{fig:validacao_literatura} apresenta a comparação dos padrões de bi-espectro obtidos com os reportados por \cite{fackrell1995}.

\begin{figure}[H]
    \centering
    \includegraphics[width=0.8\textwidth]{figures/validacao_literatura.png}
    \caption{Comparação com resultados da literatura}
    \label{fig:validacao_literatura}
\end{figure}

Os resultados mostram excelente concordância qualitativa com os padrões reportados na literatura, validando a fidelidade física das simulações realizadas.

\subsection{Análise de Sensibilidade}

A análise de sensibilidade dos parâmetros de falha revelou que:

\begin{itemize}
    \item A severidade da trinca tem impacto significativo na magnitude dos picos do bi-espectro.
    \item O desalinhamento afeta principalmente as frequências de acoplamento entre modos.
    \item A velocidade de rotação influencia a localização dos picos no bi-espectro.
\end{itemize}

\section{Classificação por Aprendizado de Máquina}

\subsection{Desempenho dos Classificadores}

A Tabela \ref{tab:desempenho_classificadores} apresenta o desempenho dos diferentes algoritmos de classificação implementados.

\begin{table}[H]
    \centering
    \caption{Desempenho dos classificadores}
    \begin{tabular}{lcccc}
        \toprule
        Algoritmo & Acurácia & Precisão & Recall & F1-Score \\
        \midrule
        SVM & 0.94 & 0.93 & 0.94 & 0.94 \\
        Random Forest & 0.91 & 0.90 & 0.91 & 0.91 \\
        Neural Network & 0.89 & 0.88 & 0.89 & 0.89 \\
        XGBoost & 0.95 & 0.94 & 0.95 & 0.95 \\
        \bottomrule
    \end{tabular}
    \label{tab:desempenho_classificadores}
\end{table}

O XGBoost apresentou o melhor desempenho geral, seguido pelo SVM. A Figura \ref{fig:matriz_confusao} apresenta a matriz de confusão para o melhor classificador.

\begin{figure}[H]
    \centering
    \includegraphics[width=0.6\textwidth]{figures/matriz_confusao.png}
    \caption{Matriz de confusão do classificador XGBoost}
    \label{fig:matriz_confusao}
\end{figure}

\subsection{Análise de Robustez}

A análise de robustez foi realizada adicionando ruído gaussiano aos dados de entrada. Os resultados mostram que o desempenho dos classificadores permanece satisfatório até níveis de ruído de 10\% do sinal original.

\section{Discussão dos Resultados}

\subsection{Contribuições Metodológicas}

Os principais resultados obtidos demonstram que:

\begin{enumerate}
    \item A simulação numérica utilizando ROSS é capaz de gerar dados de alta fidelidade para análise HOS.
    \item Os bi-espectros são eficazes para identificar assinaturas não lineares características de diferentes tipos de falhas.
    \item A metodologia proposta é robusta e pode ser aplicada a diferentes tipos de sistemas rotativos.
    \item Os classificadores de aprendizado de máquina são capazes de automatizar o processo de diagnóstico com alta precisão.
\end{enumerate}

\subsection{Limitações Identificadas}

Algumas limitações foram identificadas durante o desenvolvimento:

\begin{itemize}
    \item A modelagem de falhas complexas pode requerer refinamento adicional do modelo numérico.
    \item O tempo computacional para simulações de longa duração pode ser significativo.
    \item A validação experimental seria necessária para confirmação final dos resultados.
\end{itemize}

\subsection{Comparação com Métodos Convencionais}

A comparação com métodos convencionais de análise espectral revelou que:

\begin{itemize}
    \item Os métodos HOS são mais sensíveis a não linearidades do que a análise espectral convencional.
    \item A classificação automática oferece vantagens em termos de velocidade e consistência.
    \item A abordagem por simulação reduz a dependência de dados experimentais.
\end{itemize}

\section{Análise de Casos de Estudo}

\subsection{Caso 1: Detecção de Trinca Incipiente}

O primeiro caso de estudo focou na detecção de trincas incipientes (severidade < 5\%). Os resultados mostram que mesmo trincas de baixa severidade são detectáveis através da análise HOS, com sensibilidade superior aos métodos convencionais.

\subsection{Caso 2: Diagnóstico de Múltiplas Falhas}

O segundo caso de estudo analisou a capacidade de diagnóstico quando múltiplas falhas estão presentes simultaneamente. Os resultados indicam que a metodologia é capaz de identificar e classificar múltiplas falhas com precisão adequada.

\subsection{Caso 3: Variações Operacionais}

O terceiro caso de estudo investigou o impacto de variações nas condições operacionais (velocidade, carga, temperatura). Os resultados mostram que a metodologia é robusta a variações típicas das condições operacionais.

\section{Perspectivas de Aplicação}

\subsection{Aplicações Industriais}

A metodologia desenvolvida tem potencial para aplicação em:

\begin{itemize}
    \item Sistemas de monitoramento contínuo em plantas industriais.
    \item Desenvolvimento de sistemas de manutenção preditiva.
    \item Treinamento de algoritmos de diagnóstico com dados sintéticos.
    \item Validação de modelos numéricos antes de experimentos físicos.
\end{itemize}

\subsection{Extensões Futuras}

Possíveis extensões da metodologia incluem:

\begin{itemize}
    \item Incorporação de mais tipos de falhas (rolamentos, engrenagens).
    \item Desenvolvimento de interfaces gráficas para usuários finais.
    \item Integração com sistemas de IoT para monitoramento em tempo real.
    \item Aplicação a sistemas mais complexos (multi-rotores, acoplamentos).
\end{itemize}
